% Everything you set once for the whole thesis.

% --- Your details: used on the title page and in the PDF properties. ---
\newcommand{\myTitle}{A Clear Title That Says What the Thesis Is About}
\newcommand{\mySubtitle}{An optional subtitle, or delete it}
\newcommand{\myName}{Your Name}
\newcommand{\myDegree}{Doctor of Philosophy}
\newcommand{\myProf}{Prof.\ Dr.\ Supervisor Name}
\newcommand{\myUni}{Example University}
\newcommand{\myFaculty}{Faculty of Science}
\newcommand{\myTime}{March 2026}

% --- classicthesis options. ---
% style=classicthesis is the default look; arsclassica, linedheaders and plain are alternatives.
% dottedtoc: dotted leaders between contents entries and page numbers.
% eulerchapternumbers: chapter numbers in the Euler face instead of Palatino.
% eulermath: Euler for formulas as well (a matter of taste: remove it for plain Palatino maths).
% drafting (add it while writing): prints the compile date in the footer of every page.
\usepackage[
  style=classicthesis, dottedtoc, eulerchapternumbers, eulermath
]{classicthesis}

% PDF title and author. classicthesis has already loaded hyperref with coloured links
% (green citations, blue links); add `hidelinks` inside \hypersetup for a print version.
\hypersetup{pdftitle={\myTitle}, pdfauthor={\myName}}

% cleveref (\cref) must come after hyperref, so it is loaded here rather than in main.tex.
\usepackage[nameinlink]{cleveref}
